


\subsection[Memperkenalkan distribusi $t$]
    {Memperkenalkan distribusi $\pmb{t}$}
\label{introducingTheTDistribution}

\index{distribusi t@distribusi $t$|(}
\index{distribusi!t@$t$|(}

Dalam praktiknya, kita tidak dapat menghitung galat baku secara langsung
untuk $\bar{x}$ karena kita tidak mengetahui simpangan baku
populasi,~$\sigma$.
Kita menghadapi masalah serupa ketika menghitung galat
baku untuk proporsi sampel, yang bergantung pada proporsi
populasi,~$p$.
Solusi kita dalam konteks proporsi adalah menggunakan suatu nilai
sampel sebagai pengganti
nilai populasi ketika menghitung galat baku.
Kita akan memakai strategi serupa untuk menghitung galat
baku $\bar{x}$, dengan menggunakan simpangan
baku sampel $s$ sebagai pengganti $\sigma$:
\begin{align*}
SE = \frac{\sigma}{\sqrt{n}} \approx \frac{s}{\sqrt{n}}
\end{align*}
Strategi ini cenderung bekerja baik ketika kita memiliki
banyak data dan dapat mengestimasi $\sigma$ secara akurat menggunakan $s$.
Namun, estimasi tersebut kurang presisi pada sampel yang lebih kecil,
dan hal ini menimbulkan masalah ketika distribusi
normal digunakan untuk memodelkan $\bar{x}$.
% --
%when the sample size is large --
%but it is less reliable when the sample size is smaller
%than about 30. % independent observations.

Kita akan mendapati bahwa distribusi baru yang disebut
\termsub{distribusi $\pmb{t}$}{distribusi t@distribusi $t$} berguna
untuk perhitungan inferensi.
Distribusi~$t$, yang ditampilkan sebagai garis penuh dalam
Gambar~\ref{tDistCompareToNormalDist}, berbentuk lonceng.
Namun, ekornya lebih tebal daripada ekor distribusi normal,
yang berarti observasi lebih mungkin jatuh lebih dari dua
simpangan baku dari rata-rata dibandingkan pada distribusi
normal. %\footnote{The standard deviation of the
  %$t$-distribution is actually a little more than 1.
  %However, it is useful to always think of the $t$-distribution
  %as having a standard deviation of 1 in all of our applications.}
%This distribution is important since it accounts for
%a key challenge with modeling the sample mean:
%the standard error of the sample mean isn't as
%precise when the sample size is small.
Ketebalan ekstra pada ekor distribusi $t$ merupakan
koreksi yang diperlukan untuk mengatasi masalah penggunaan~$s$
sebagai pengganti $\sigma$ dalam perhitungan $SE$.

\begin{figure}[h]
  \centering
  \Figure[Ditampilkan distribusi normal baku dan distribusi t. Distribusi t juga berbentuk lonceng, tetapi puncaknya lebih rendah dan lebih landai daripada distribusi normal serta ekornya lebih tebal. Sebagai contoh, terdapat bagian distribusi yang cukup besar -- mungkin 5\% untuk distribusi t khusus ini -- yang membentang hingga di bawah -3 dan di atas 3, sedangkan distribusi normal sangat mendekati nol pada nilai di bawah -3 atau di atas 3.]{0.7}{tDistCompareToNormalDist}
  \caption{Perbandingan distribusi $t$
      dan distribusi normal.}
  \label{tDistCompareToNormalDist}
\end{figure}

Distribusi $t$ selalu berpusat di nol dan
memiliki satu parameter: derajat kebebasan.
\termsub{Derajat kebebasan ($\pmb{df}$)}
    {derajat kebebasan ($df$)!distribusi $t$}
menentukan bentuk tepat distribusi $t$ yang berbentuk lonceng.
Beberapa distribusi $t$ ditampilkan dalam
Gambar~\ref{tDistConvergeToNormalDist}
sebagai perbandingan dengan distribusi normal.

Secara umum, kita akan menggunakan distribusi $t$
dengan $df = n - 1$ untuk memodelkan rata-rata sampel
ketika ukuran sampelnya $n$.
Artinya, ketika kita memiliki lebih banyak observasi,
derajat kebebasannya akan lebih besar dan
distribusi $t$ akan semakin menyerupai
distribusi normal baku;
ketika banyaknya derajat kebebasan sekitar 30 atau lebih,
distribusi $t$ hampir tidak dapat dibedakan
dari distribusi normal.

\begin{figure}[h]
  \centering
  \Figure[Empat distribusi t dengan derajat kebebasan 1, 2, 4, dan 8 ditampilkan bersama distribusi normal dalam grafik yang sama. Semakin besar derajat kebebasan, semakin dekat distribusi t dengan distribusi normal; puncak di pusat semakin tinggi dan mendekati puncak distribusi normal, sedangkan ekor distribusinya tampak semakin tipis.]{0.75}{tDistConvergeToNormalDist}
  \caption{Semakin besar derajat kebebasan,
      distribusi $t$ semakin menyerupai distribusi
      normal baku.}
  \label{tDistConvergeToNormalDist}
\end{figure}

\begin{onebox}{Derajat kebebasan
    ($\pmb{\MakeLowercase{df}}$)}
  Derajat kebebasan menentukan bentuk
  distribusi $t$.
  Semakin besar derajat kebebasan, semakin baik
  distribusi tersebut menghampiri model normal. \stdvspace{}

  Ketika memodelkan $\bar{x}$ menggunakan distribusi $t$,
  gunakan $df = n - 1$.
\end{onebox}

%\Comment{Cut this next sentence?}
%In Section~\ref{tDistSolutionToSEProblem},
%we relate degrees of freedom to sample size.

Distribusi $t$ memberi kita fleksibilitas lebih besar daripada
distribusi normal ketika menganalisis data numerik.
Dalam~praktiknya, perangkat lunak statistik lazim digunakan,
seperti R, Python, atau SAS untuk analisis ini.
Sebagai alternatif, kalkulator grafik atau
\termsub{tabel $\pmb{t}$}{tabel t@tabel $t$} dapat digunakan;
tabel $t$ serupa dengan tabel distribusi normal,
dan dapat ditemukan dalam Lampiran~\ref{tDistributionTable},
yang memuat petunjuk penggunaan dan contoh
bagi mereka yang ingin menggunakan opsi ini.
Apa pun pendekatan yang dipilih, terapkan metode Anda
pada contoh-contoh di bawah ini untuk memastikan
pemahaman Anda mengenai distribusi $t$.

\begin{examplewrap}
\begin{nexample}{Berapa proporsi distribusi $t$
    dengan 18 derajat kebebasan yang berada di bawah -2.10?}
  Seperti dalam soal peluang normal, pertama-tama kita menggambar
  grafik dalam Gambar~\ref{tDistDF18LeftTail2Point10}
  dan mengarsir daerah di bawah -2.10.
%  If this were a normal distribution, the area would be
%  a little less than 0.025, since about 5\% of the area
%  under a normal curve goes out beyond $\pm 1.96$ standard
%  deviations.
  Dengan perangkat lunak statistik, kita dapat memperoleh
  nilai yang presisi: 0.0250.
%  The tail area below -2.10 in the $t$-distribution with
%  $df = 18$ is the same as the tail area below -1.96 in
%  the normal distribution.
\end{nexample}
\end{examplewrap}

\begin{figure}
  \centering
  \Figure[Ditampilkan distribusi t dengan 18 derajat kebebasan; daerah di bawah -2.10 diarsir dan tampak mewakili sekitar 2\% hingga 5\% dari luas distribusi. Secara umum, ketika derajat kebebasan lebih besar daripada sekitar 10, seperti dalam kasus ini, perbedaan visual antara distribusi t dan distribusi normal sangat halus, meskipun perbedaan tersebut tetap penting untuk perhitungan kita.]{0.42}{tDistDF18LeftTail2Point10}
  \caption{Distribusi $t$ dengan 18 derajat kebebasan.
      Daerah di bawah -2.10 telah diarsir.}
  \label{tDistDF18LeftTail2Point10}
\end{figure}

\begin{examplewrap}
\begin{nexample}{Distribusi $t$ dengan 20 derajat kebebasan
    ditampilkan pada panel kiri
    Gambar~\ref{tDistDF20RightTail1Point65}.
    Estimasikan proporsi distribusi yang berada
    di atas 1.65.}
  Pada distribusi normal, nilainya akan sekitar
  sekitar~0.05, sehingga kita mengharapkan distribusi $t$
  memberikan nilai di sekitar itu.
  Dengan perangkat lunak statistik: 0.0573.
\end{nexample}
\end{examplewrap}

\begin{figure}
  \centering
  \Figure[Dua distribusi t ditampilkan pada dua grafik terpisah. Grafik pertama menampilkan distribusi t dengan 20 derajat kebebasan dan daerah di atas 1.65 diarsir; daerah itu tampak mewakili sekitar 5\% dari luas total distribusi. Grafik kedua menampilkan distribusi t dengan 2 derajat kebebasan dan daerah di bawah -3 serta di atas 3 diarsir. Karena derajat kebebasannya sangat kecil, ekor distribusi ini jauh lebih tebal dan pusatnya juga berpuncak lebih tajam. Masing-masing ekor tersebut tampak mewakili sekitar 2\% hingga 5\% dari luas di bawah distribusi ini.]{0.72}{tDistDF20RightTail1Point65}
  \caption{Kiri: Distribusi $t$ dengan 20 derajat
      kebebasan, dengan daerah di atas 1.65 diarsir.
      Kanan:~Distribusi $t$ dengan 2 derajat kebebasan,
      dengan daerah yang berjarak lebih dari 3 satuan dari 0 diarsir.}
  \label{tDistDF20RightTail1Point65}
\end{figure}

\begin{examplewrap}
\begin{nexample}{Distribusi $t$ dengan 2 derajat kebebasan
    ditampilkan pada panel kanan
    Gambar~\ref{tDistDF20RightTail1Point65}.
    Estimasikan proporsi distribusi yang berjarak lebih
    dari 3~satuan dari rata-rata (di atas atau di bawah).}
  Dengan derajat kebebasan yang begitu sedikit, distribusi $t$
  akan menghasilkan nilai yang jauh berbeda dari hasil distribusi
  normal.
  Pada distribusi normal, luasnya akan sekitar
  0.003 berdasarkan aturan 68-95-99.7.
  Untuk distribusi $t$ dengan $df = 2$, luas pada kedua
  ekor di luar 3~satuan berjumlah 0.0955.
  Luas ini sangat berbeda dari luas yang
  kita peroleh dari distribusi normal.
\end{nexample}
\end{examplewrap}

\begin{exercisewrap}
\begin{nexercise}
Berapa proporsi distribusi $t$ dengan 19 derajat
kebebasan yang berada di atas -1.79 satuan?
Gunakan metode pilihan Anda untuk mencari luas ekor.\footnotemark{}
\end{nexercise}
\end{exercisewrap}
\footnotetext{Kita ingin mencari daerah yang diarsir di \emph{atas}
  -1.79 (pembuatan gambarnya kami serahkan kepada Anda).
  Luas ekor bawah adalah 0.0447,
  sehingga luas daerah atas adalah $1 - 0.0447 = 0.9553$.}

\index{distribusi!t@$t$|)}
\index{distribusi t@distribusi $t$|)}


%\subsection{Conditions for using the $\mathbf{t}$-distribution
%    for inference on a sample mean}
%\label{tDistSolutionToSEProblem}
%
%\noindent%
%To proceed with the $t$-distribution for inference about a single mean, we first check two conditions.
%\begin{description}
%\item[Independence.]
%    We verify this condition just as we did before.
%    We collect a simple random sample, or if the data are from
%    an experiment or random process, we check to the best of our
%    abilities that the observations were independent.
%\item[Sample size.]
%    We use the earlier rule of thumb to evaluate this condition:
%
%    If the sample size $n$ is less than 30
%    and there are no clear outliers in the data,
%    then the sample size condition is satisfied.
%
%    If the sample size $n$ is at least 30
%    and there are no \emph{particularly extreme} outliers,
%    then the sample size condition is satisfied.
%\end{description}
%When examining a sample mean and estimated standard error
%from a sample of $n$ independent and nearly normal observations,
%we use a $t$-distribution with $n - 1$ degrees of freedom~($df$).
%For example, if the sample size was 19, then we would use the
%$t$-distribution with $df = 19 - 1 = 18$ degrees of freedom
%and proceed in a way similar to how we worked with proportions.
